\newcommand{\ModuleID}{EL-III.4}
\newcommand{\ModuleTitle}{Secrets and Postcommunions: Gift and Effect}
\newcommand{\ModuleStage}{Stage III — Independent Missal Reading}
\newcommand{\ModuleUnit}{III.4}
\newcommand{\ModuleOutcome}{Distinguish offering, sign, reception, present effect, and promised fruition through syntax and lexical families}
\newcommand{\ModulePrerequisites}{EL-III.3 and reliable Collect architecture}
\newcommand{\ModuleExtensionPractice}{Worksheets III.14 and III.15 remain optional reception/fruition and solemnity practice}
\newcommand{\ModuleNext}{EL-III.5, Lessons and Epistles}
\newcommand{\ModuleMastery}{On an unseen Secret or Postcommunion, map the petition and its effects, resolve pronouns from morphology, and distinguish what the prayer offers, receives, signifies, or seeks without replacing grammar with theological inference.}
\newcommand{\ModuleDelayNote}{}
\newcommand{\ModuleLessonSource}{\CourseRoot/shared/content/volume3/all}
\newcommand{\ModuleExerciseSource}{\CourseRoot/shared/exercises/volume3_4/volume3}
\newcommand{\ModuleWorkedBank}{III}
\newcommand{\ModuleExerciseBank}{III}
\newcommand{\ModuleWorksheetVariants}{A,D}
\newcommand{\ModuleScopeNote}{This is grammatical and lexical genre study, not a sacramental or theological exposition.}
\newcommand{\ModulePractice}{%
  \SubmoduleExtension{The offered sign}
    {The noun that signifies and the reality signified may occupy different syntactic slots.}
    {On the lesson example, label finite verb, subject, objects, antecedents, and every term assigned to gift, sign, or reality.}
    {State which link is grammatical and which is a semantic classification.}
    {The model proves the parse from endings, agreement, and government, then adds semantic-role labels. A semantic association alone may not supply an antecedent or case.}
  \SubmoduleExtension{Reception as a sign of future fruition}
    {Postcommunion prose often relates reception now to a sought effect without collapsing their time or role.}
    {Make a two-line now/future map for the lesson's Postcommunion: put each finite or non-finite form on the line supported by its morphology and governor.}
    {Identify one English tense choice that could blur the Latin relation.}
    {A passing model distinguishes completed or present reception from the petition's desired fruition and cites the decisive forms. Natural English may vary, but it may not reverse that relation.}
  \SubmoduleExtension{A future feast approached now}
    {An ambiguous neuter form must be resolved by government and clause structure before symbolism is discussed.}
    {List every formal possibility for the lesson's ambiguous form, cross out those excluded by its governor, and explain the surviving relation.}
    {Why is thematic expectation weaker evidence than government?}
    {The model begins with morphology, then identifies the governing word and required case or function. Theme can confirm a grammatical result but cannot override it.}
  \SubmoduleExtension{Semantic families}
    {A family ledger prevents one English gloss from flattening distinct acts of offering, receiving, cleansing, and enjoying.}
    {Create four family entries from the lesson with lemma, printed form, construction, agent or participant, and local sense.}
    {Give one pair that should remain separate despite overlapping English translations.}
    {Every accepted entry cites a form and construction. The contrast must explain a difference in agency, object, temporal orientation, or genre role rather than merely list synonyms.}
}
